\documentclass[12pt,a4paper,english,twoside]{article}
\usepackage{fontspec}
\usepackage[english]{babel}
\usepackage{amsmath}
\usepackage{amsfonts}
\usepackage{amssymb}
\usepackage{graphicx}
\usepackage{fancyhdr}
\usepackage{lmodern}
\usepackage{minted}
\usepackage{hyperref}
\usepackage{verbatim}
\usepackage{alltt}
\usepackage{color}
\usepackage{listings}
\lstset{frame=single,keepspaces=true,numbers=none,tabsize=2}
\defaultfontfeatures{Ligatures=TeX}
\setmainfont{DejaVu Serif}
\setsansfont{DejaVu Sans}
\setmonofont{DejaVu Sans Mono}
\usepackage[cm]{fullpage}
\setlength{\headsep}{10pt}
\setlength{\footskip}{20pt}
\addtolength{\textheight}{-30pt}
\author{
    Bart de Boer \texttt{bart.de.boer@hightechict.nl}
    \and
    Bas Bossink \texttt{bas.bossink@hightechict.nl}
}
\date{\today}
\title{Functional Programming and Test Driven Development\\
\vspace{2 mm} {\large A match made in heaven}}
\newminted{scheme}{linenos=true,mathescape,xleftmargin=20px}
\usepackage{natbib}  \bibliographystyle{chicago}
\renewcommand{\headrulewidth}{0.5pt}
\renewcommand{\headheight}{35pt}
\fancyhead[L]{\tiny Functional Programming and Test Driven Development}
\fancyhead[R]{\includegraphics[scale=0.2]{HTLogo}}
\fancyfoot[C]{\thepage}
\pagestyle{fancy}
\begin{document}
\maketitle
\section{Scheme the language}
This section will briefly describe a small portion of the Scheme language that 
will be used in this workshop.
Before we start a few typographical conventions all text written in a 
\texttt{monospaced font} is either a Scheme keyword or text provided or 
returned from the Scheme interpreters REPL (Read Eval Print Loop), a console 
application that evaluates Scheme expressions as you provide them. The prompt 
of the REPL starts with a hash \texttt{\#}, followed by a semicolon \texttt{;}, the 
character that starts a comment in Scheme, followed by a counter of the number 
of successfully evaluated expressions, followed by the greater than \texttt{>} 
symbol. For example \texttt{\#;37>}.
\subsection{Introduction}
Since Scheme is a LISP (LISt Processing) dialect the central syntactical 
element in the language is the list. But before we dive into lists lets look 
at Scheme's `primitive' expressions, that is expression that evaluate to 
themselves. In Scheme these are: booleans, numbers, characters, strings, 
symbols and vectors.
\subsection{Boolean values}
Scheme has the boolean values \texttt{\#f} and \texttt{\#t} for false and true 
respectively. In boolean expressions all values except \texttt{\#f} are 
treated as true.
\subsection{Numbers}
Scheme supports the following types of numbers
\begin{itemize}
\item integers with base 2, 8, 10, 16
\begin{lstlisting}
#;5>#b100101
37
#;6>#o45
37
#;7>37
37
#;8>#x25
37
\end{lstlisting}
\item rationals
\begin{lstlisting}
#;9>23/5
23/5
#;10>#x20/1a
16/13
\end{lstlisting}
\item floating point values
\begin{lstlisting}
#;11>1.3552e-12
1.3552e-12 
#;12>.2342
0.2342
\end{lstlisting}
\item complex numbers
\begin{lstlisting}
#;13>1+37i
1+37i
#;14>37@1.7
-4.76724628893441+36.6915979867413i
\end{lstlisting}
\end{itemize}
\subsection{Characters}
Character literals are preceded by the 
\texttt{\#\textbackslash} characters.
\begin{lstlisting}
#;15>#\a
#\a
#;16>#\newline
#\newline
#;17>#\Z
#\Z
\end{lstlisting}
\subsection{Strings}
Strings are written as a number of characters surrounded by double-quotes.
\begin{lstlisting}
#;18>"Fred Flinstone"
"Fred Flinstone"    
#;19>"Barney Rubbles"
"Barney Rubbles"
\end{lstlisting}
\subsection{Symbols}
A symbol is an immutable string for which strings with the same content 
will refer to the same object. Symbols can be compared with 
\texttt{eq?} (essentially pointer comparison), whereas string need to be 
compared by looking at their contents with \texttt{equal?}. Let's look at some 
examples:
\begin{lstlisting}
#;20>'Fred
Fred
#;21>(string->symbol "Fred Flinstone")
|Fred Flinstone|
#;22>(eq? 'Fred 'Fred)
#t
#;23>(eq? "Fred" "Fred")
#f
#;24>(equal? "Fred" "Fred")
#t
\end{lstlisting}
\subsection{Other types of objects}
\subsubsection{Lists}
A list is denoted as a \texttt{'('} followed by elements separated by spaces 
followed by a \texttt{')'}. Since the list is also the primary syntactical 
construct of the Scheme language lists are interpreted in a special way. The 
interpreter tries to evaluate the first element of the list to a procedure and 
applies that procedure with the rest of the elements in the list as the 
arguments to the procedure.
\begin{lstlisting}
#;22>(1 2 3)
Error: call of non-procedure: 1

Call history:

<syntax>	  (1 2 3)
<eval>	  (1 2 3)	<--

\end{lstlisting}
In the example output you can see the way Scheme interprets lists. The first 
element of the list is interpreted as a procedure if it is not one of the 
standard syntactic forms. This interpretation can be circumvented by 
\texttt{quote}-ing the expression using the \texttt{quote} syntactic form or 
it's abbreviated form \texttt{'}. An example of this is given below:
\begin{lstlisting}
#;22>(quote (1 2 3))
(1 2 3)
\end{lstlisting}
\begin{lstlisting}
#;23>'(1 2 3)
(1 2 3)
\end{lstlisting}
\begin{lstlisting}
#;24>(+ 1 2 3)
6
\end{lstlisting}
\paragraph{Creating lists}
Lists can be created in several different ways:
\begin{itemize}
\item List literals \begin{lstlisting}
#;25>(quote (1 2 3))
(1 2 3)
#;26>'(1 2 3)
(1 2 3)
\end{lstlisting}
\item Using the \texttt{list} procedure
\begin{lstlisting}
#;27>(list 1 2 3)
(1 2 3)
\end{lstlisting}
\item Using the \texttt{cons} procedure to construct lists  \begin{lstlisting}
#;28>(cons 1 (cons 2 (cons 3 '())))
(1 2 3)
\end{lstlisting}
The name \texttt{cons} stands for "construct". The \texttt{cons} procedure 
creates what is known as a cons cell, a pair of pointers. These pointer pairs 
can be used to create a single linked list where the second part of the pair 
points to the next element (= cons cell) in a list. The last cons cell of a 
list points to a distinguished value that is not a pair in our case the empty 
list \texttt{()}.  \end{itemize}
\paragraph{Deconstructing lists}
Lists can be deconstructed using the \texttt{car} and \texttt{cdr} procedures. 
These names stem from the original Lisp implementation in the late 1950's.  
\begin{itemize}
\item \texttt{car} short for "Contents of the Address part of Register number" 
  is referred to as \texttt{head} or \texttt{first} in some languages.
\item \texttt{cdr} short for "Contents of the Decrement part of Register 
  number" is referred to as \texttt{tail} or \texttt{rest} in some languages.
\end{itemize} \begin{lstlisting}
#;29>(car '(1 2 3))
1
#;30>(cdr '(1 2 3))
(2 3)
\end{lstlisting}
\subsubsection{Functions}
In Scheme procedures or functions are first class values as well we'll look at 
them in detail later but first we are going to look at definitions.
\subsection{Bindings a.k.a Definitions}
\subsubsection{Top-level Bindings}
To create a top-level binding the \texttt{define} syntactic form is used which 
in it's most basic form looks like this:
\begin{equation*}
\texttt{(define <variable> <expression>)}
\end{equation*}
Some examples:
\begin{lstlisting}
#;31>(define fred 37)
#;32>fred
37
#;33>(define barney "Betty")
#;34>barney
"Betty"
\end{lstlisting}

\subsubsection{Local Bindings}
Local bindings can be introduced using the \texttt{let}, \texttt{let*} and 
\texttt{letrec} syntactical forms each of them will be described below.
\begin{equation*}
\texttt{(let ((name$_{0}$ <expression$_{0}$>) $\dots$ (name$_{n}$ 
<expression$_{n}$>)) <body>)}
\end{equation*}
\begin{lstlisting}
#;35>(let ((a 4) (b 3)) (* a b))
12
#;36>(let ((a 4) (b (* a a))) (* a b)) 
Error: unbound variable: a

    Call history:

    <syntax>	  (let ((a 4) (b (* a a))) (* a b))
    <syntax>	  (##core#let ((a 4) (b (* a a))) (* a b))
    <syntax>	  (##core#begin (* a b))
    <syntax>	  (* a b)
    <syntax>	  (* a a)
    <eval>	  (* a a)	<--
#;37>(let* ((a 4) (b (* a 3))) (* a b))
48
#;38>(letrec
  ((odd? (lambda (x) (if (= x 0) #f (even? (- x 1)))))
   (even? (lambda (x) (if (= x 0) #t (odd? (- x 1)))))) (odd? 37))
#t
\end{lstlisting}
As you can see from the interpreter output in a \texttt{let} binding the 
different bindings can not refer to each other. In a \texttt{let*} however the 
bindings are evaluated from left to right and the previous bindings are in the 
scope of the subsequent bindings. The \texttt{letrec} binding can be used to 
define mutually recursive bindings as shown in the example above. The value of 
a \texttt{let} binding is value of the last expression in the body. 

\subsection{Procedures/Functions}
\subsubsection{Creating procedures}
Procedures can be created using the \texttt{lambda} syntactical form. Here 
lambda refers back to the lambda calculus \cite{lambda}.
The basic syntax of a lambda expression looks like this:
\begin{equation*}
\texttt{(lamdba (par$_{0} \dots$ par$_{n}$) <body>)}
\end{equation*}
Examples of \texttt{lambda} expressions:
\begin{lstlisting}
#;38>(lambda (x) (* x x))
#<procedure (? x)>
#;39>((lambda (x) (* x x)) 2)
4 
#;40>((lambda (x y) (* x y)) 3 4)
12
#;41>(define square (lambda (x) (* x x)))
#;43>(square 3)
9
\end{lstlisting}
From the last example you can take away that combining \texttt{define} with a 
\texttt{lambda} expression gives you a way to define functions. Similarly you 
could use a \texttt{let} binding with a \texttt{lambda} expression to define a 
`local' function.
\subsubsection{Applying procedures}
We have already seen examples of applying a procedure to arguments above. 
There also is a handy function called \texttt{apply} that can be used to apply 
a procedure to it's arguments.
\begin{lstlisting}
#;44>(apply (lambda (x y) (* x y)) '(3 4))
12
#;45>((lambda (x y) (* x y)) (car '(3 4)) (car (cdr '(3 4))))
12
\end{lstlisting} Using \texttt{apply} is sometimes handy because then you 
don't have to deconstruct the list as you can see in the example above.
\section{Fizzbuzz}
\subsection{Writing a solution for the basic Fizzbuzz kata}
Now that we have covered the most basic forms of Scheme we will look at some 
more Scheme syntactic forms using the Fizzbuzz kata as a vehicle. The result 
of completing this kata will be used in your next assignment.
\subsubsection{The assignment}
Write a function that for each positive integer $x$ will produce the following 
result:
\begin{enumerate}
  \item if $x$ is divisible by 3 return the string "fizz"
  \item if $x$ is divisible by 5 return the string "buzz"
  \item if $x$ is divisible by both 3 and 5 return the string "fizzbuzz"
  \item otherwise return $x$
\end{enumerate}
\subsection{Our first test}
Since this workshop is about Test Driven Development we will start with our 
first test. We will be using a Scheme `standard library' called SRFI 78. It is 
already provided for us in the check.scm file. A first test would look like 
this:
\subsubsection{A first test case}
\begin{schemecode}
(load "check.scm")
(check (fizzbuzz 1) => 1)
(check-report)
\end{schemecode}
Saving this in a file called fizzbuzz.scm We can run this test using a 
prepared script \texttt{run-tests} like so:
\begin{lstlisting}
>run-tests fizzbuzz.scm
\end{lstlisting}
This will result in the following output:
\begin{lstlisting}

Error: unbound variable: fizzbuzz

    Call history:

    <syntax>	  (##core#lambda () (fizzbuzz 1))
    <syntax>	  (##core#begin (fizzbuzz 1))
    <syntax>	  (fizzbuzz 1)
    <syntax>	  (##core#undefined)
    <eval>	  (>=426 check:mode427 1)
    <eval>	  (check:proc428 (quote429 (fizzbuzz 1) )...
    <eval>	  [check:proc] (eqv?60 tmp57 (quote 0))
    <eval>	  [check:proc] (eqv?60 tmp57 (quote 1))
    <eval>	  [check:proc] (eqv?60 tmp57 (quote 10))
    <eval>	  [check:proc] (eqv?60 tmp57 (quote 100))
    <eval>	  [check:proc] (check:report-expression expression)
    <eval>	  [check:report-expression] (newline)
    <eval>	  [check:report-expression] (check:write expression)
    <eval>	  [check:report-expression] (display " => ")
    <eval>	  [check:proc] (thunk)
    <eval>	  (fizzbuzz 1)	<--
\end{lstlisting}
Our test can not be interpreted because we have not yet defined our 
\texttt{fizzbuzz} procedure. Lets look at the code and explain each line 
one by one:

\begin{enumerate}
  \item line 1: loads our unit testing library of choice, the check library
  \item line 2: this is our first assert you could read this as: verify that calling 
    the \texttt{fizzbuzz} function with \texttt{1} will result in the value 
    \texttt{1}
 \item line 3: this last line creates a report of our results.
\end{enumerate}
\subsubsection{The first fizzbuzz implementation}
Now we have to define the first revision of our \texttt{fizzbuzz} function, 
to keep things simple we will write our production code in the same file as 
our test code. Here is the first implementation of our \texttt{fizzbuzz} 
function:
\begin{schemecode}
(load "check.scm")

(define fizzbuzz (lambda (x) 1))

(check (fizzbuzz 1) => 1)

(check-report)
\end{schemecode}
\begin{lstlisting}
>run-tests fizzbuzz.scm

(fizzbuzz 1) => 1 ; correct

; *** checks *** : 1 correct, 0 failed.
\end{lstlisting}
\subsubsection{The second test case}
Great our first passing test. Lets add our next test case.  
\begin{schemecode}
(load "check.scm")

(define fizzbuzz (lambda (x) 1))

(check (fizzbuzz 1) => 1)
(check (fizzbuzz 2) => 2)

(check-report)
\end{schemecode}
\begin{lstlisting}
>run-tests fizzbuzz.scm

(fizzbuzz 1) => 1 ; correct

(fizzbuzz 2) => 1 ; *** failed ***
 ; expected result: 2

; *** checks *** : 1 correct, 1 failed. First failed example:

(fizzbuzz 2) => 1 ; *** failed ***
 ; expected result: 2

\end{lstlisting}
\subsubsection{Making our second test case pass}
To save some space in this document we will use the \texttt{report-failed} 
mode of our test library in the next runs. Making this failing test pass is a 
matter of replacing our constant with a parameter.
\begin{schemecode}
(load "check.scm")
(check-set-mode! 'report-failed)

(define fizzbuzz (lambda (x) x))

(check (fizzbuzz 1) => 1)
(check (fizzbuzz 2) => 2)

(check-report)
\end{schemecode}
\begin{lstlisting}
>run-tests fizzbuzz.scm

; *** checks *** : 2 correct, 0 failed.  
\end{lstlisting}
Excellent. Now lets add our next test case.
\subsubsection{The third test case}
\begin{schemecode}
(load "check.scm")
(check-set-mode! 'report-failed)

(define fizzbuzz (lambda (x) x))

(check (fizzbuzz 1) => 1)
(check (fizzbuzz 2) => 2)
(check (fizzbuzz 3) => "fizz")

(check-report)
\end{schemecode}
\begin{lstlisting}
>run-tests fizzbuzz.scm

(fizzbuzz 3) => 3 ; *** failed ***
 ; expected result: "fizz"

; *** checks *** : 2 correct, 1 failed. First failed example:

(fizzbuzz 3) => 3 ; *** failed ***
 ; expected result: "fizz"
\end{lstlisting}

To be able to make the failing test pass we need to know how we can write 
conditional expressions in Scheme, luckilly that is the topic of the next 
section.

\subsection{Conditionals}
\paragraph{\texttt{if}}
The simplest conditional expression in Scheme is the \texttt{if} expression 
which
has the following syntax:
\begin{equation*}
    \texttt{(if <expression> <then expression> [<else expression>])}
\end{equation*}
The value of the if expression is the value of the expression that gets 
evaluated depending on the value of the \texttt{<expression>}, the else 
expression is optional. If the \texttt{<expression>} evaluates to false and 
the else expression is missing than the value of the if expression is 
undefined.

Now we can make our failing test pass:
\subsubsection{Making the third test case pass}
\begin{schemecode}
(load "check.scm")
(check-set-mode! 'report-failed)

(define fizzbuzz (lambda (x) (if (= 3 x) "fizz" x)))

(check (fizzbuzz 1) => 1)
(check (fizzbuzz 2) => 2)
(check (fizzbuzz 3) => "fizz")

(check-report)
\end{schemecode}
\begin{lstlisting}
>run-tests fizzbuzz.scm

; *** checks *** : 3 correct, 0 failed.
\end{lstlisting}
The next interesting case is of course the number 5.  
\subsubsection{The fourth test case}
\begin{schemecode}
(load "check.scm")
(check-set-mode! 'report-failed)

(define fizzbuzz (lambda (x) (if (= 3 x) "fizz" x)))

(check (fizzbuzz 1) => 1)
(check (fizzbuzz 2) => 2)
(check (fizzbuzz 3) => "fizz")
(check (fizzbuzz 5) => "buzz")

(check-report)
\end{schemecode}
To make it pass we could use a nested if but Scheme provides us with the 
\texttt{cond} expression which reads a lot nicer.
\paragraph{\texttt{cond}}
Below the syntax for the \texttt{cond} syntactic form:
\begin{lstlisting}[mathescape]
  (cond ((<predicate>$_{0}$) <expression>$_{0}$) $\vdots$ ((<predicate$_{n}$) 
  <expression>$_{n}$) [(else <expression$_{n+1}$)])
\end{lstlisting}
Each of the predicates gets evaluated in order of appearance the value of the 
\texttt{cond} expression is the value of the expression of the first predicate 
that evaluated to true. A \texttt{cond} expression can have an optional 
\texttt{else} clause that will be evaluated if none of the predicates 
evaluated to true.
\subsubsection{Making the fourth test case pass}
\begin{schemecode}
(load "check.scm")
(check-set-mode! 'report-failed)

(define fizzbuzz (lambda (x) (cond ((= 3 x) "fizz")
          ((= 5 x) "buzz")
          (else x))))

(check (fizzbuzz 1) => 1)
(check (fizzbuzz 2) => 2)
(check (fizzbuzz 3) => "fizz")
(check (fizzbuzz 5) => "buzz")

(check-report)
\end{schemecode}
\begin{lstlisting}
>run-tests fizzbuzz.scm

; *** checks *** : 4 correct, 0 failed.  \end{lstlisting}
The rest of the kata can now be implemented, all the Scheme constructs that we 
will need have been introduced. All that we need to know is, does Scheme have 
something like a `remainder' or `modulus' function. The function to use here 
is the \texttt{modulo} function. Below you find the finished implementation:
\subsubsection{The completed solution for our FizzBuzz kata}
\begin{schemecode}
(load "check.scm")
(check-set-mode! 'report-failed)

(define fizzbuzz 
  (lambda (x) 
    (define divisible-by?  
      (lambda (divisor number)
        (= 0 (modulo number divisor))))
    (cond ((divisible-by? 15 x) "fizzbuzz")
        ((divisible-by? 3 x) "fizz")
        ((divisible-by? 5 x) "buzz")
          (else x))))

(check (fizzbuzz 1) => 1)
(check (fizzbuzz 2) => 2)
(check (fizzbuzz 3) => "fizz")
(check (fizzbuzz 5) => "buzz")
(check (fizzbuzz 6) => "fizz")
(check (fizzbuzz 7) => 7)
(check (fizzbuzz 10) => "buzz")
(check (fizzbuzz 15) => "fizzbuzz")
(check (fizzbuzz 30) => "fizzbuzz")
(check-report)
\end{schemecode}
\section{Assignment l: removing code duplication}
\subsection{Running tests}
To make running tests easy during this workshop we will be using a 
\texttt{watchr} script that will automatically run our tests when they are 
saved. When you start up the virtual machine you will be greated with a split 
console window. The top asks you which editor you would like to use during 
this session, the lower part will contain the output of running the our tests.
\subsection{DRY-ing up test code using higher order functions}
During this excercise we will remove some duplication (DRY= Don't Repeat 
Yourself) from the test code we created earlier using higher order functions 
(functions that take other functions as parameters or return funtions as their 
result). The code for this exercise is in the \texttt{fizzbuzz.scm} file.  
Start by opening it using any of the available editors from the menu.  
\subsection{`Extract method'}
Remove some of the duplication by extracting a function that takes as 
parameters the input integer and the expected output. This function should 
invoke the test method such that calling this new function executes a single 
test case. Use this function in each of the test cases.
\subsection{Reduce duplication further using \texttt{map}}
Now we will remove even more duplication by using the \texttt{map} and 
\texttt{apply} functions. We saw the \texttt{apply} function in the 
introduction it takes a function and a list as arguments and applies the 
supplied function to the supplied arguments in the list. The \texttt{map} 
function takes a function and a list as its parameters and returns a new list 
consisting of the function applied to each of the elements in the list.
Here is a short example
\begin{lstlisting}
#;46>(map (lambda (x) (* x x)) '(1 2 3))
(1 4 9)
\end{lstlisting}
The assignment is to use the \texttt{map} and \texttt{apply} functions to 
rewrite the tests as an applycition of \texttt{map} on a list of test cases.
You should use a \texttt{let} binding to avoid poluting the global namespace.
\section{Assignment 2: Legacy code}
Some of you may know the legendary Hewlett-Packard calculators from the good 
old days that use Reverse Polish Notation (RPN). A form of notation that first 
denotes the operands followed by the operator. This assignment works with a 
RPN calculator
written in Scheme. It uses a stack to hold intermediary results, it can read a 
Scheme expression pushes all numbers it encounters onto the stack and when it 
encounters an operator pops the 2 operands of the stack applies the operator 
and pushes the result back onto the stack. Below you find some examples of 
using this calculator.
\begin{lstlisting}
#;47>(load "stack-calculator.scm")
#;48>(rpn '(37 42 23 + *))
2405
\end{lstlisting}
Lets first briefly look over the code:
\begin{schemecode}
(define (rpn expression)
  (define operator? 
    (lambda (expression)
      (let ((token (car expression)))
        (or (eqv? '+ token)
            (eqv? '- token)
            (eqv? '* token)
            (eqv? '/ token)))))
  (define operator 
    (lambda (expression) 
      (car expression)))
  (define calculate 
    (lambda (exp stack)
      (if (null? expression)
        "there is nothing to calculate!"
        (cond ((null? exp) (car stack))
              ((operator? exp)
               (let* ((first-operand (cadr stack))
                      (second-operand (car stack))
                      (result (eval (list 
                                      (operator exp) 
                                      first-operand 
                                      second-operand))))
                 (calculate (cdr exp) (cons result (cddr stack)))))
              (else (calculate (cdr exp) (cons (car exp) stack)))))))
  (calculate expression '()))
\end{schemecode}
This code contains a few functions that we have not yet discussed:
\begin{enumerate}
  \item lines 5-8: \texttt{eqv?} tests the equivalency of its argument since here in lines 
    5-8 only symbols are used this means equality
  \item line 14.16: \texttt{null?} returns true if its parameter is the empty list
  \item line 18: \texttt{cadr} is a composition of \texttt{cdr} and \texttt{car} 
    \begin{equation*}
      \texttt{(cadr alist)} \Rightarrow \texttt{(car (cdr list))}
    \end{equation*}
  \item line 20: \texttt{eval} evaluates the given list in the same manner as the 
    Scheme interpreter would, which means evaluate the first element of the 
    list to a procedure definition evaluate the rest of the elements and 
    supply them as the formal parameters to the procedure.
  \item line 24: \texttt{cddr} is a composition of \texttt{cdr} with itself
    \begin{equation*}
      \texttt{(cddr alist)} \Rightarrow \texttt{(cdr (cdr alist))}
    \end{equation*}
\end{enumerate}
\subsection{Building the safety net}
Start by writing some tests that validate the current behavior of the 
\texttt{rpn} function. When you feel you have captured enough of the current 
behavior of the function proceed to the next assignment.
\subsection{Adding the \texttt{\^{}} operator}
Extend the \texttt{rpn} function such that it can understand expresions that 
contain the \texttt{\^{}} operator which should calculate the power function.  
\begin{equation*}
\texttt{(rpn '(2 3 \^{}))} \Rightarrow 2^{3} \Rightarrow 8 
\end{equation*}

\subsubsection{Baby steps}
To accomplish the above task it can be helpful if you promote the 
\texttt{operator?} and \texttt{operator} functions to top level bindings such 
that you can write separate test cases for these individual functions.
\subsection{Adding the \texttt{sqrt} function}
Upto now the calculator only understands binary operators, extend the 
calculator such that expressions can contain the \texttt{sqrt} operator that 
applies the \texttt{sqrt} function to the operand on the top of the stack.  
\begin{equation*}
  \texttt{(rpn '(9 sqrt))} \Rightarrow \sqrt{9} \Rightarrow 3.0
\end{equation*}
The Scheme function that calculates the square root of its argument is called 
\texttt{sqrt}.
\begin{enumerate}
  \item Create the \texttt{unary-operator?} function that can tell if the 
    first element in the expression is a unary operator or not. 
  \item Use this function in the \texttt{operator?} function.
  \item Extend the \texttt{operator} function to return the \texttt{sqrt} function as a 
    translation for \texttt{sqrt}.
  \item Extend the \texttt{calculate} function to take unary operators into 
    account.
  \item Introduce an evaluate function that calls \texttt{eval}. You need to 
    make this a function with a variable number of arguments. Such a function 
    can be defined using the follwing syntax
    \begin{equation*}
      \texttt{(lambda (<parameter>$_{0} \dots $<parameter>$_{n}$ . <rest>) <body>)}
    \end{equation*}
    In the body of the \texttt{lambda} expression the \texttt{rest} parameter 
    will be bound to a list containing the parameters$_{n+1}$ and beyond.
\end{enumerate}
\bibliography{handout}
\end{document}
